\section{MXNet：深度学习框架的百花齐放}

\subsection{从 CXXNet 到 MXNet}

AlexNet 的成功重新点燃了陈天奇对深度学习的热情。他和许斌（现 NVIDIA）开发了 CXXNet，使用 C++ 模板编程（expression template）技术生成内联 GPU kernel。这个技术的核心思想是：你在写 update rule（比如做乘法、减法等运算）时，它可以通过 C++ 模板把这些操作压缩成一个 kernel，生成内联的 GPU 代码。CXXNet 的性能不错，还被用来打 Kaggle 深度学习相关的比赛。

与此同时，深度学习框架领域正处于``百花齐放''的阶段：

\begin{itemize}
    \item \textbf{张真 + 敏捷}（MSR/NYU）：开发了 Minerva，Python first 的设计理念在当时非常先进
    \item \textbf{李沐}（CMU）：做 Parameter Server（参数服务器），需要分布式作为 first class citizen
    \item \textbf{张池元}（MIT）：用 Julia 写了一个叫 Moka 的框架
    \item \textbf{林敏}（NUS）：做了 Purine 框架
    \item \textbf{杨青}（UC Berkeley）：Caffe 已经 release，是当时的主流框架，但核心是 C++
\end{itemize}

通过 NeurIPS 等学术会议的契机，这些分散在 UW、CMU、MIT、NUS、NYU 的 PhD 学生聚在一起，决定合力做一个统一的框架------这就是 MXNet 的起源。

\begin{knowledgebox}{MXNet 的名字由来}
MXNet 是 Minerva 和 CXXNet 的混合（Mix），同时也代表 ``mix programming style''------混合声明式计算图（declarative，先声明计算图再执行）和命令式编程（imperative，写了就直接算）。当时这两种范式存在明显的 trade-off：声明式对系统优化有利（可以做全局优化），命令式对开发者更友好（Debug 简单）。MXNet 的选择是两者都支持。后来 PyTorch 选择了纯命令式（eager mode），TensorFlow 1.x 选择了纯声明式（后来 2.0 也转向了 eager mode），MXNet 一开始就提供了混合模式。
\end{knowledgebox}

陈天奇还提到，UW 的同组同学中有 Joe Redmon（YOLO 的创始人），``他在博士之前就做了 YOLO''。AlexNet 出来后，组里很多人都开始探索深度学习方向。

\subsection{核心设计理念}

MXNet 的三大创新：
\begin{enumerate}
    \item \textbf{Python First}：在 Caffe 以 C++ 为核心的年代，MXNet 选择以 Python 为一等公民。
    \item \textbf{自动调度}：根据计算依赖关系自动实现多卡并行，通信和计算可以 overlap。
    \item \textbf{自动求导}：从 Caffe 的手写 backpropagation 进化到基于计算图的自动求导。
\end{enumerate}

\subsection{与 TensorFlow 的竞争}

陈天奇暑假在 Google Brain 实习时，发现 TensorFlow 即将发布。他是 TensorFlow 的 dogfood 用户之一，同时也向 Brain 团队分享了 MXNet 的设计经验。但团队并没有退缩：``出生牛犊不怕虎，我们也可以 compete 一下。''甚至``很有激情地继续往前推动''。

时间线上，MXNet 的立项早于 TensorFlow 的发布。大家在前一年的 NeurIPS 上聚在一起决定往前走，TensorFlow 在 2015 年底 release，MXNet 的正式发布约在同期。最终 MXNet 只投了一篇 workshop paper，作者按字母序排列------``纯粹就是觉得一起做一个项目挺好的''。

\begin{knowledgebox}{PhD 学生的``纯粹''合作}
主持人问：这么多 PhD 聚在一起，没有人管事，怎么能 work out？陈天奇回答：``根本没有人可以说了算。因为大家纯粹是想要一起做项目。''这种基于共同兴趣、没有层级关系的合作模式，在学术开源社区中反而有其优势------每个人都是出于技术热情而非 KPI 在贡献。合作中``就技术讲技术其实很简单，直接讨论方案然后往前走就可以了''。
\end{knowledgebox}

\begin{warningbox}{MXNet 淡出历史舞台的教训}
MXNet 最终被 PyTorch 超越，陈天奇总结了两个关键教训：
\begin{enumerate}
    \item \textbf{用户体验优先于性能}：MXNet 在多卡并行上长期是最优秀的框架（NVIDIA 跑 MLPerf 都用它），但在用户体验上不如 PyTorch。``在选择用户体验和性能上面，我们觉得都要，但其实应该先选用户体验。''
    \item \textbf{工程人力不足}：为支持每一代 GPU 的迭代，需要大量工程投入，PhD 学生团队难以持续。这直接促使他转向 TVM------\textbf{用编译技术来减少工程开销}。
\end{enumerate}
\end{warningbox}

\subsection{李沐的中流砥柱作用}

陈天奇特别提到李沐是 MXNet 项目中做了``非常大的努力和牺牲''的人。``其实最后李沐是在这里面推动最大的。''李沐毕业后去了 Amazon，推动 MXNet 成为 AWS 官方框架，并开发了 GluonCV 和 GluonNLP。陈天奇指出 GluonNLP ``in some sense 有点像 Hugging Face''------提供各种预训练模型的统一接口。但最终这些生态层面的努力没能改变 PyTorch 社区已经壮大的格局。

一个不太为人知的事实是：NVIDIA 在非常长的一段时间里都在 contribute MXNet、使用 MXNet。原因是在所有开源框架中，MXNet 的多卡并行性能一直是最优秀的。NVIDIA 在做 MLPerf 跑分提交时，``一定要拿到最好分数的时候，一直在用 MXNet''------直到 ResNet 不再是 benchmark 的主角之后才逐渐淡出。

陈天奇还提到图森未来的王乃岩和侯晓迪``有很长一段时间一直在帮忙贡献和支持 MXNet''。这个项目的成功和最终淡出，都离不开社区中很多人的贡献。

\subsection{本章小结}

MXNet 的经历教会陈天奇两件事：（1）用户体验比技术指标更重要；（2）工程人力的瓶颈需要通过更根本的技术手段（编译器）来解决，而不是简单地堆人。这两个教训直接催生了 TVM。

